\section{Representasi State}

Pada permasalahan ini, state yang digunakan untuk mewakili kebijakan suku bunga BI akan direpresentasikan dalam bentuk vektor suku bunga acuan BI-Rate untuk seluruh horizon kebijakan sekaligus yang dinyatakan oleh

\[
    s = (r_1, r_2, \dots, r_T).
\]

Tiap elemen $r_t$ adalah BI-Rate yang berlaku pada kuartal ke-$t$, dengan $T = 8$ periode (2 tahun). Domain nilai $r_t$ didiskritisasi dalam kelipatan 25 basis poin (0,25\%), sesuai dengan konvensi dalam RDG BI yang umumnya mengubah suku bunga dalam kelipatan tersebut. Secara matematis, haruslah nilai $r_t$ memenuhi

\[
    r_t \in \{3{,}00\%,\ 3{,}25\%,\ 3{,}50\%,\ \dots,\ 8{,}00\%\}
\]

Ukuran ruang state dari formulasi ini adalah $21^{8} \approx 3,78 \times 10^{10}$ kemungkinan (21 tingkat suku bunga diskrit, 8 periode). Ukuran ini, tentunya akan sangat tidak efisien jika ditelusuri jika menggunakan algoritma \textit{brute-force}. Dengan demikian, diperlukan algoritma \textit{local search} untuk mengeksplorasi ruang pencarian secara lebih strategis.

\section{Aturan dan Batasan Permasalahan}

Permasalahan ini dibatasi oleh lima \textit{constraint} yang merepresentasikan kondisi neraca keuangan riil Indonesia, mencakup aspek domestik maupun eksternal. \textit{Constraint} C1 bersifat struktural terhadap representasi state itu sendiri sehingga dapat diverifikasi langsung dari vektor $s$ tanpa perlu simulasi. \textit{Constraint} C2 sampai C5 melibatkan variabel makroekonomi hasil simulasi sehingga pelanggarannya baru dapat dideteksi setelah fungsi \texttt{compute\_objective(s)} dijalankan.

\textbf{C1. Batas dan Diskritisasi Suku Bunga.}

\[
    r_t \in [r_{min}, r_{max}] = [3{,}00\%,\ 8{,}00\%], \quad r_t \bmod 0{,}25\% = 0 \qquad \forall t
\]

\textit{Constraint} ini memastikan setiap $r_t$ berada pada rentang suku bunga yang realistis untuk ekonomi berkembang dan selalu dalam kelipatan 25 basis poin sesuai konvensi perubahan suku bunga BI.

\textbf{C2. Batas Defisit Transaksi Berjalan.}

\[
    CA_t(s) \geq -\theta, \qquad \theta = 3{,}00\%
\]

Defisit transaksi berjalan tidak boleh melebihi ambang 3,00\% PDB, merujuk pada batas aman IMF untuk negara \textit{emerging market}. \textit{Constraint} ini merepresentasikan kondisi terkini Indonesia yang baru mengalami defisit transaksi berjalan pertama setelah 72 bulan surplus berturut-turut, sehingga ruang gerak kebijakan yang memperlebar defisit lebih jauh perlu dibatasi.

\textbf{C3. Batas \textit{Spread} Imbal Hasil Obligasi Pemerintah.}

\[
    y_t^{SBN} - r_t \leq \sigma_{max}, \qquad \sigma_{max} = 2{,}50\%
\]

Selisih antara imbal hasil SBN dan BI-Rate tidak boleh melebar melebihi $\sigma_{max} = 2{,}50$ poin persentase, untuk menghindari risiko \textit{fiscal dominance}, yaitu situasi di mana biaya pembiayaan utang pemerintah membengkak akibat sentimen pasar terhadap arah kebijakan moneter. Risiko ini secara eksplisit menjadi perhatian dalam diskusi perluasan kewenangan BI membeli SBN di pasar perdana pada UU P2SK.

\textbf{C4. Konvergensi Target Inflasi di Akhir Horizon.}

\[
    \left| \pi_T - \pi^{*} \right| \leq 1\%, \qquad \pi^{*} = 2{,}5\%
\]

Lintasan kebijakan harus membawa inflasi kembali ke dalam rentang sasaran resmi $2{,}5\% \pm 1\%$ pada akhir horizon $T = 8$. \textit{Constraint} ini menjamin bahwa state yang dihasilkan bukan sekadar optimal secara matematis, tetapi juga kredibel sebagai kebijakan publik sesuai kesepakatan sasaran inflasi Pemerintah-BI untuk periode 2026-2027.

\textbf{C5. Gradualisme Kebijakan (\textit{Penalized Interest Rate Smoothing}).}

\[
    \left| r_t - r_{t-1} \right| \leq 50 \text{ bps} \qquad \forall t
\]

Berbeda dengan C1-C4 yang bersifat \textit{hard constraint} (state harus memenuhi agar dinyatakan \textit{feasible}), C5 diperlakukan sebagai \textit{penalized constraint}, dimana perubahan suku bunga melebihi 50 bps diizinkan, namun dikenakan penalti kuadratik pada fungsi objektif. Pendekatan ini mencerminkan realitas kebijakan moneter, yaitu perubahan suku bunga yang terlalu agresif dapat mengguncang ekspektasi pasar dan merusak kredibilitas \textit{forward guidance} \cite{woodford2003}, namun dalam kondisi luar biasa (seperti krisis) bank sentral tetap dapat mengambil langkah drastis.

\subsection{Definisi \textit{Feasible} dan Peran $\pi_T$}

\textbf{Status \textit{feasible}.} Sebuah state $s$ dinyatakan \textbf{\textit{feasible}} jika dan hanya jika seluruh \textit{hard constraint} C1, C2, C3, dan C4 terpenuhi. \textit{Constraint} C5 (gradualisme) dikeluarkan dari definisi \textit{feasible} karena bersifat \textit{soft/penalized}, dimana perubahan suku bunga lebih dari 50 bps tidak membatalkan kelayakan kebijakan, melainkan hanya dikenakan penalti pada fungsi objektif $J_{penalized}(s)$. Secara formal,

\[
    \text{feasible}(s) \iff \text{C1} \land \text{C2} \land \text{C3} \land \text{C4}
\]

\textbf{Peran $\pi_T$ (inflasi terminal).} $\pi_T$ adalah nilai inflasi pada akhir horizon kebijakan ($t = T = 8$). Nilai ini menjadi variabel kunci karena

\begin{itemize}
    \item $\pi_T$ digunakan untuk verifikasi \textit{constraint} C4, yaitu $|\pi_T - \pi^{*}| \leq 1\%$. Constraint ini menjamin bahwa lintasan suku bunga yang dihasilkan algoritma tidak hanya optimal secara matematis, tetapi juga kredibel sebagai kebijakan publik. Inflasi harus kembali ke dalam rentang sasaran $2{,}5\% \pm 1\%$ pada akhir horizon.
    \item $\pi_T$ dilaporkan bersama hasil akhir setiap algoritma agar pembaca dapat menilai secara langsung apakah kebijakan yang direkomendasikan berhasil mengarahkan inflasi menuju target.
    \item Bersama RMSD$_{\pi}$, $\pi_T$ memberikan gambaran lengkap. RMSD$_{\pi}$ akan mengukur deviasi rata-rata sepanjang seluruh horizon, sedangkan $\pi_T$ menunjukkan kondisi akhir yang dicapai.
\end{itemize}

\section{\textit{Succesor Function}}

\textit{Successor function} didefinisikan sebagai cara membangkitkan state tetangga dari sebuah state $s$. Dalam kerangka \textit{local search}, tetangga adalah state yang ``dekat'' dengan state saat ini menurut suatu metrik tertentu. Pada permasalahan ini, akan digunakan tiga jenis \textit{move}.
\begin{itemize}
    \item \textit{Move 1. Perubahan Titik.} Operator ini mengubah satu elemen $r_t$ sebesar $\pm 25$ basis poin (satu tingkat diskritisasi) dan mempertahankan elemen lain tetap. Karena dinamika model bersifat kausal, perubahan pada $r_t$ mempengaruhi seluruh \textit{trajectory} dari periode $t$ hingga $T$. Secara implementatif, setelah $r_t$ dimodifikasi, simulasi ekonomi dijalankan ulang dari periode $t$ sampai $T$. \textit{Move} ini bersifat lokal dan presisi, cocok untuk \textit{fine-tuning} di sekitar solusi yang sudah cukup baik.
    \item \textit{Move 2. Pergeseran \textit{Window}.} Operator ini memilih satu periode $t$ dan menggeser seluruh sub-state dari $r_t$ hingga $r_T$ sebesar $\pm 25$ basis poin. \textit{Move} ini mensimulasikan perubahan arah kebijakan secara menyeluruh. Menggeser ke atas merepresentasikan pergeseran ke arah yang lebih \textit{hawkish} (mengetatkan kebijakan), sedangkan menggeser ke bawah merepresentasikan pergeseran ke arah yang lebih \textit{dovish} (melonggarkan kebijakan). Karena pergeseran bersifat seragam pada seluruh sisa horizon, \textit{move} ini mempertahankan \textit{constraint} C2 selama pergeseran tidak menyebabkan pelanggaran batas atas atau bawah.
    \item \textit{Move 3. Crossover.} Operator ini khusus digunakan pada \textit{Genetic Algorithm} dan bekerja pada dua state induk. Sebuah titik potong $k \in \{1, \dots, T-1\}$ dipilih, lalu segmen $[r_1, \dots, r_k]$ dari induk pertama ditukar dengan segmen yang sama dari induk kedua, menghasilkan dua state anak. \textit{Move} ini memungkinkan pertukaran blok kebijakan antar kandidat, mempercepat eksplorasi dengan menggabungkan segmen yang baik dari dua state berbeda.
\end{itemize}



Untuk algoritma \textit{Hill-Climbing} dan \textit{Simulated Annealing}, tetangga dibangkitkan dengan memilih secara acak antara \textit{Move 1} dan \textit{Move 2} pada tiap iterasi. Untuk \textit{Genetic Algorithm}, populasi terdiri dari sekumpulan state $s$, dengan mutasi menggunakan \textit{Move 1} dan \textit{crossover} menggunakan \textit{Move 3}. Perlu dicatat bahwa seluruh \textit{move} harus tetap menjaga \textit{constraint} C1 setelah modifikasi. Jika hasil modifikasi melanggar C1, state tetangga diperbaiki terlebih dahulu (dengan \textit{clipping} ke rentang yang diizinkan) sebelum dievaluasi.

\section{\textit{Objective Function}}

\subsection{Definisi \textit{Objective Function}}

Fungsi objektif $J(s)$ adalah ukuran kuantitatif yang menentukan seberapa baik sebuah state $s$ dalam menyeimbangkan tiga tujuan kebijakan moneter Bank Indonesia, yaitu stabilitas harga, stabilitas nilai tukar, dan daya beli masyarakat. \textit{Local search} bertugas mencari state yang memaksimalkan $J(s)$ pada ruang state $\mathcal{S}$ yang memenuhi seluruh \textit{constraint} C1 sampai C5.

\[
    s^{*} = \arg\max_{s \,\in\, \mathcal{S}} J(s)
\]

Nilai $J(s)$ dihitung melalui dua langkah. Pertama, state $s$ disimulasikan pada model makroekonomi untuk menghasilkan \textit{trajectory} variabel $\{y_t, \pi_t, e_t, PP_t, CA_t, y_t^{SBN}\}_{t=1}^{T}$. Kedua, \textit{trajectory} tersebut dievaluasi terhadap target kebijakan dan \textit{constraint} untuk menghasilkan vektor suku bunga akhir.

\subsection{\textit{Loss Function} Teoretik}

\textit{Loss function} $L(s)$ mengikuti kerangka \textit{quadratic loss function} Svensson \cite{svensson1997} yang menjadi standar analisis kebijakan bank sentral. Secara konseptual, $L(s)$ mengukur total ``kerugian'' atau deviasi variabel makroekonomi dari targetnya sepanjang horizon kebijakan. Setiap komponen dihitung sebagai kuadrat deviasi, sehingga deviasi besar dikenakan penalti yang jauh lebih berat daripada deviasi kecil (sifat kuadratik). Perumusan ini mencerminkan preferensi bank sentral yang lebih memilih deviasi kecil di banyak periode dibandingkan satu deviasi besar di satu periode.

\[
    L(s) = \sum_{t=1}^{T} \delta^{t} \Big[\, w_\pi (\pi_t - \pi^{*})^2 + w_y \, y_t^{2} + w_{pp} (PP_t - PP^{*})^2 + w_r (r_t - r_{t-1})^2 \,\Big]
\]

Keempat suku di dalam kurung siku memiliki interpretasi kebijakan yang spesifik. Suku $w_\pi (\pi_t - \pi^{*})^2$ menghukum deviasi inflasi dari sasaran $\pi^{*} = 2{,}5\%$. Suku $w_y \, y_t^{2}$ menghukum \textit{output gap} yang tidak nol pada kedua arah, karena ekonomi yang \textit{overheating} ($y_t > 0$) maupun resesi ($y_t < 0$) sama-sama tidak diinginkan. Suku $w_{pp} (PP_t - PP^{*})^2$ menghukum deviasi daya beli dari target $PP^{*} = 0$, yaitu kondisi ideal di mana pertumbuhan riil tepat mengimbangi inflasi sehingga daya beli masyarakat tidak tergerus. Suku $w_r (r_t - r_{t-1})^2$ menghukum volatilitas suku bunga yang berlebihan, mendorong \textit{interest rate smoothing} sebagaimana dikehendaki \textit{constraint} C2 secara lunak pada fungsi objektif.

Faktor diskonto $\delta = 0{,}95$ memberi bobot lebih besar pada periode awal dibanding periode akhir. Dengan $\delta < 1$, deviasi pada kuartal pertama memiliki kontribusi lebih signifikan terhadap total \textit{loss} dibanding deviasi yang sama pada kuartal kedelapan, mencerminkan preferensi bank sentral terhadap stabilitas jangka pendek dan ketidakpastian yang meningkat seiring bertambahnya horizon proyeksi.

Bobot $w_\pi, w_y, w_{pp}, w_r$ menentukan prioritas relatif antar tujuan kebijakan. Konfigurasi default $(w_\pi, w_y, w_{pp}, w_r) = (1{,}0,\ 0{,}5,\ 0{,}3,\ 0{,}2)$ menempatkan stabilitas inflasi sebagai prioritas tertinggi, diikuti oleh \textit{output gap}, daya beli, dan \textit{smoothing} suku bunga sebagai prioritas terendah.

\subsection{\textit{Loss Function} Final}

\textit{Loss function} final yang digunakan dalam implementasi adalah $J_{penalized}(s)$, yaitu negatif dari $L(s)$ yang telah ditambah penalti \textit{constraint}. Negasi dilakukan karena algoritma \textit{local search} secara konvensional bekerja dengan memaksimalkan fungsi objektif, sedangkan $L(s)$ adalah fungsi yang ingin diminimalkan.

\[
    J_{penalized}(s) = -L(s) - \mu \sum_{i=2}^{5} \max\!\big(0,\; \text{violation}_i(s)\big)^2
\]

Suku pertama, $-L(s)$, adalah negatif dari \textit{quadratic loss} yang mengubah permasalahan minimisasi $L(s)$ menjadi maksimisasi. Seluruh suku pada $L(s)$ bernilai non-negatif karena berbentuk kuadrat, sehingga $J(s) = -L(s) \leq 0$ untuk semua $s$. Nilai $J(s)$ yang semakin mendekati nol menandakan \textit{trajectory} ekonomi yang semakin dekat ke target. Nilai maksimum teoretis $J(s) = 0$ hanya tercapai jika semua deviasi sama dengan nol di seluruh periode, kondisi yang secara praktis mustahil karena kehadiran guncangan stokastik $\epsilon_t$.

Suku kedua adalah penalti untuk \textit{constraint} C2 sampai C5 yang dilanggar. Pendekatan \textit{penalty method} dipilih karena \textit{constraint} C2 sampai C5 melibatkan variabel hasil simulasi sehingga tidak dapat diverifikasi sebelum simulasi dijalankan, berbeda dengan C1 yang dapat dicek langsung dari vektor $s$. Bobot penalti $\mu$ ditetapkan sebesar 100, cukup besar untuk mendominasi perbaikan pada $L(s)$ sehingga algoritma terdorong kuat menuju region \textit{feasible} sebelum melakukan \textit{fine-tuning} terhadap nilai $J(s)$. State yang memenuhi seluruh \textit{constraint} memiliki suku penalti nol sehingga $J_{penalized}(s) = -L(s)$. Setiap \textit{constraint} $i \in \{2, 3, 4, 5\}$ memiliki pelanggaran $\text{violation}_i(s)$ yang dihitung sebagai berikut.
\begin{align*}
    \text{violation}_2(s) & = \max_{t}\!\big(0,\ -\theta - CA_t\big)                   \\
    \text{violation}_3(s) & = \max_{t}\!\big(0,\ (y_t^{SBN} - r_t) - \sigma_{max}\big) \\
    \text{violation}_4(s) & = \max\!\big(0,\ |\pi_T - \pi^{*}| - 1\%\big)              \\
    \text{violation}_5(s) & = \max_{t}\!\big(0,\ |r_t - r_{t-1}| - 0{,}50\%\big)
\end{align*}

Pelanggaran dikuadratkan agar pelanggaran besar dikenakan penalti yang jauh lebih berat daripada pelanggaran kecil, mendorong algoritma untuk memprioritaskan perbaikan pada \textit{constraint} yang paling parah dilanggar.


\section{Formulasi \textit{Initial State}}

Initial state $s_0$ dibentuk secara acak dari seluruh rentang yang diizinkan. Setiap $r_t$ dipilih secara acak dari distribusi uniform pada rentang $[r_{min}, r_{max}]$ lalu dibulatkan ke kelipatan 0,25\% terdekat, sehingga \textit{constraint} C1 secara implisit terpenuhi. Jika hasil pembangkitan acak tidak memenuhi C2 sampai C4, state tetap diterima sebagai initial state dalam status \textit{infeasible}, dan proses \textit{local search} bertugas mengarahkan pencarian menuju state yang \textit{feasible} sekaligus optimal melalui mekanisme penalti pada fungsi objektif.

\medskip
\noindent\textbf{Kondisi awal sistem (periode $t = 0$).}
\begin{align*}
    r_0      & = 5{,}75\%                                                                                        &  & \text{BI-Rate}                   \\
    \pi_0    & = 3{,}34\%                                                                                        &  & \text{Inflasi tahunan Juni 2026} \\
    y_0      & = 0{,}00\%                                                                                        &  & \text{\textit{Output gap} awal}  \\
    e_0      & = 17{.}885\ \text{Rp/USD}                                                                         &  & \text{Nilai Dollar ke Rupiah}    \\
    r_t^{US} & = (3{,}375\%,\ 3{,}125\%,\ 3{,}125\%,\ 3{,}125\%,\ 3{,}375\%,\ 3{,}625\%,\ 3{,}875\%,\ 4{,}125\%) &  & \text{Fed Rate}
\end{align*}

\medskip
\noindent\textbf{Parameter constraint dan target.}
\begin{align*}
    T            & = 8        &  & \text{Horizon kebijakan (kuartalan, 2 tahun)}        \\
    r_{min}      & = 3{,}00\% &  & \text{Batas bawah suku bunga}                        \\
    r_{max}      & = 8{,}00\% &  & \text{Batas atas suku bunga}                         \\
    \pi^{*}      & = 2{,}5\%  &  & \text{Sasaran inflasi tengah}                        \\
    \Delta_{min} & = 1{,}50\% &  & \text{Selisih minimum BI-Rate terhadap Fed Rate}     \\
    \theta       & = 3{,}00\% &  & \text{Batas defisit transaksi berjalan (\% PDB)}     \\
    \sigma_{max} & = 2{,}50\% &  & \text{\textit{Spread} maksimum SBN terhadap BI-Rate} \\
    \delta       & = 0{,}95   &  & \text{Faktor diskonto \textit{loss function}}
\end{align*}

\medskip
\noindent\textbf{Bobot komponen fungsi objektif.}
\begin{align*}
    w_{\pi} & = 1{,}0 &  & \text{Bobot inflasi}                          \\
    w_{y}   & = 0{,}5 &  & \text{Bobot \textit{output gap}}              \\
    w_{pp}  & = 0{,}3 &  & \text{Bobot daya beli}                        \\
    w_{r}   & = 0{,}2 &  & \text{Bobot \textit{interest rate smoothing}}
\end{align*}

\section{\textit{Model Evaluation Metric}}

Kualitas solusi akhir tidak dapat diukur dari $J_{penalized}(s)$ saja, karena nilai $J(s)$ dipengaruhi oleh bobot $w$ dan $\delta$ yang meskipun bersifat tetap, merepresentasikan satu konfigurasi preferensi spesifik. Diperlukan metrik evaluasi yang berdiri di luar fungsi objektif, bersifat netral, dan langsung mengukur seberapa dekat \textit{trajectory} inflasi terhadap target kebijakan. Metrik tersebut adalah \textbf{\textit{Root Mean Square Deviation} (RMSD)}.

Sebelum membahas RMSD, perlu diklarifikasi pemisahan antara komponen yang bersifat tetap (\textit{fixed}) dan komponen yang bebas diatur (\textit{free}) dalam eksperimen.

\begin{itemize}
    \item \textit{Komponen tetap.} Seluruh parameter yang tertuang dalam Formulasi Initial State pada spesifikasi bersifat wajib dan tidak boleh diubah. Ini mencakup parameter struktural model ($\kappa, \beta, \phi, \rho_y, \rho_\pi, \rho_e, \alpha_1, \alpha_2, \rho_{fiskal}$), kondisi awal ($r_0, \pi_0, y_0, e_0, r_0^{US}$), target ($\pi^{*}, PP^{*}$), parameter \textit{constraint} ($r_{min}, r_{max}, \Delta_{min}, \theta, \sigma_{max}$), bobot fungsi objektif ($w_\pi, w_y, w_{pp}, w_r$), faktor diskonto ($\delta$), horizon ($T$), dan protokol \textit{shock} (\texttt{seed = 42}). Lingkungan simulasi yang identik menjamin perbedaan nilai RMSD semata-mata mencerminkan perbedaan kualitas algoritma, bukan perbedaan ``dunia'' yang disimulasikan.
    \item \textit{Komponen bebas.} Hal-hal yang dapat divariasikan hanyalah aspek yang berkaitan dengan algoritma \textit{local search} itu sendiri, yaitu pilihan algoritma (Hill-Climbing, Simulated Annealing, atau Genetic Algorithm), jumlah iterasi (\texttt{max\_iter}), jumlah generasi dan ukuran populasi untuk GA, suhu awal dan laju pendinginan untuk SA, probabilitas mutasi, strategi seleksi, serta \textit{random seed} untuk inisialisasi state dan pemilihan \textit{neighbor}. Variasi pada komponen bebas inilah yang menghasilkan nilai $J(s)$ dan RMSD yang berbeda antar eksperimen meskipun ekonomi yang disimulasikan identik.
\end{itemize}
RMSD$_{\pi}$ mengukur akar rata-rata kuadrat deviasi inflasi dari target $\pi^{*} = 2{,}5\%$ sepanjang horizon $T = 8$. Metrik ini tidak melibatkan bobot, diskonto, maupun penalti \textit{constraint}, sehingga murni mengukur seberapa dekat \textit{trajectory} inflasi terhadap target.
\[
    \text{RMSD}_{\pi} = \sqrt{\frac{1}{T} \sum_{t=1}^{T} (\pi_t - \pi^{*})^{2}}
\]
Secara implementatif, RMSD$_{\pi}$ dihitung setelah \texttt{simulate\_economy} menghasilkan \textit{trajectory}.
\[
    \texttt{rmsd\_pi = np.sqrt(np.mean((traj["pi"] - pi\_star)**2))}
\]

Nilai RMSD$_{\pi}$ bersifat non-negatif, dengan nol sebagai nilai terbaik (inflasi tepat di target sepanjang seluruh horizon). Semakin kecil RMSD$_{\pi}$, semakin baik state suku bunga yang ditemukan. Berdasarkan karakteristik numerik dari parameter dan \textit{shock} yang ditetapkan, interpretasi kualitatif nilai RMSD$_{\pi}$ adalah sebagai berikut.

\begin{itemize}
    \item $\text{RMSD}_{\pi} < 0{,}50$ dinilai sangat baik. Inflasi rata-rata hanya menyimpang kurang dari setengah poin persentase dari target, menunjukkan algoritma \textit{local search} berhasil menemukan state yang efektif menavigasi \textit{trade-off} kebijakan.
    \item $0{,}50 \leq \text{RMSD}_{\pi} < 1{,}00$ dinilai baik dan masih dalam rentang toleransi sasaran inflasi $\pm 1\%$.
    \item $\text{RMSD}_{\pi} \geq 1{,}00$ dinilai buruk karena rata-rata deviasi inflasi melampaui batas toleransi resmi, mengindikasikan algoritma terjebak pada solusi yang tidak memenuhi \textit{constraint} C4 atau gagal mengeksplorasi ruang state secara memadai.
\end{itemize}

Selain RMSD$_{\pi}$, dapat pula dihitung RMSD$_{y}$ terhadap target $y^{*} = 0$ dan RMSD$_{PP}$ terhadap target $PP^{*} = 0$ untuk analisis lebih lanjut. Namun dalam konteks perbandingan utama, hanya RMSD$_{\pi}$ yang digunakan.

\section{Daftar Command Utama}

Berikut adalah daftar fungsi yang perlu diimplementasikan dalam program, dikelompokkan berdasarkan tiga fase utama dalam pipeline eksekusi. Setiap command ditulis dalam format \texttt{nama\_fungsi(param1, param2, ...)} diikuti oleh tujuan dan daftar parameternya.

\subsection*{Fase 1. Persiapan dan Simulasi Ekonomi}

\texttt{generate\_shocks(T=8, seed=42)}

Menghasilkan guncangan stokastik untuk seluruh horizon simulasi. Seluruh peserta wajib menggunakan \textit{seed} yang sama agar hasil dapat dibandingkan.

\begin{itemize}
    \item \texttt{T} (int) jumlah periode horizon kebijakan.
    \item \texttt{seed} (int) \textit{seed} untuk reproduktibilitas, ditetapkan 42.
    \item \textit{Return} \texttt{dict} berisi lima array \texttt{NumPy} berdimensi \texttt{(T,)} untuk guncangan IS, PC, FX, CA, dan SBN.
\end{itemize}

\texttt{simulate\_economy(s, shocks, params)}

Menjalankan simulasi model makroekonomi secara rekursif dari $t = 1$ sampai $T$. Fungsi ini adalah inti komputasi program karena dipanggil setiap kali sebuah state dievaluasi.

\begin{itemize}
    \item \texttt{s} (list[float]) state suku bunga $s = (r_1, \dots, r_T)$, tiap elemen dalam persen.
    \item \texttt{shocks} (dict) guncangan stokastik hasil \texttt{generate\_shocks}.
    \item \texttt{params} (dict) parameter model ($\kappa, \beta, \phi, \rho_y, \rho_\pi, \rho_e, \alpha_1, \alpha_2, \rho_{fiskal}$) dan kondisi awal ($\pi_0, y_0, e_0, r_0^{US}$).
    \item \textit{Return} \texttt{dict} berisi enam array berdimensi \texttt{(T,)} untuk \textit{trajectory} variabel $\{y_t, \pi_t, e_t, PP_t, CA_t, y_t^{SBN}\}_{t=1}^{T}$.
\end{itemize}

\subsection*{Fase 2. Evaluasi State}

\texttt{compute\_purchasing\_power(y, pi)}

Menghitung proksi daya beli masyarakat sebagai selisih \textit{output gap} dan inflasi.

\begin{itemize}
    \item \texttt{y} (array[float]) \textit{output gap} tiap periode.
    \item \texttt{pi} (array[float]) inflasi tiap periode.
    \item \textit{Return} \texttt{array[float]} berisi $PP_t = y_t - \pi_t$ untuk $t = 1, \dots, T$.
\end{itemize}

\texttt{compute\_current\_account(de, y, eps\_ca)}

Menghitung proksi neraca transaksi berjalan berdasarkan depresiasi nilai tukar dan \textit{output gap}.

\begin{itemize}
    \item \texttt{de} (array[float]) persen depresiasi $\Delta e_t$ tiap periode.
    \item \texttt{y} (array[float]) \textit{output gap} tiap periode.
    \item \texttt{eps\_ca} (array[float]) guncangan stokastik $\epsilon_t^{CA}$.
    \item \textit{Return} \texttt{array[float]} berisi $CA_t = \alpha_1 \Delta e_t - \alpha_2 y_t + \epsilon_t^{CA}$, dalam satuan persen PDB.
\end{itemize}

\texttt{compute\_sbn\_yield(r, eps\_sbn, rho\_fiskal)}

Menghitung proksi imbal hasil Surat Berharga Negara.

\begin{itemize}
    \item \texttt{r} (array[float]) BI-Rate tiap periode.
    \item \texttt{eps\_sbn} (array[float]) guncangan stokastik $\epsilon_t^{SBN}$.
    \item \texttt{rho\_fiskal} (float) premi risiko fiskal, ditetapkan $1{,}50$.
    \item \textit{Return} \texttt{array[float]} berisi $y_t^{SBN} = r_t + \rho_{fiskal} + \epsilon_t^{SBN}$.
\end{itemize}

\texttt{check\_constraints(s, trajectory, params)}

Memvalidasi seluruh \textit{constraint} C1 sampai C5 terhadap sebuah state dan \textit{trajectory} simulasinya, serta menghitung besar pelanggaran untuk penalti fungsi objektif. Status \textit{feasible} hanya mempertimbangkan \textit{hard constraint} C1 sampai C4; C5 adalah \textit{soft constraint} yang tidak mempengaruhi status \textit{feasible}.

\begin{itemize}
    \item \texttt{s} (list[float]) state suku bunga.
    \item \texttt{trajectory} (dict) hasil simulasi dari \texttt{simulate\_economy}.
    \item \texttt{params} (dict) parameter \textit{constraint} ($r_{min}, r_{max}, r_0, \theta, \sigma_{max}, \pi^{*}$).
    \item \textit{Return} \texttt{dict} berisi \texttt{feasible} (bool) dan \texttt{violations} (dict) dengan kunci \texttt{C2, C3, C4, C5} yang bernilai besar pelanggaran masing-masing. C1 dapat diverifikasi langsung dari \texttt{s} tanpa \texttt{trajectory}.
\end{itemize}

\texttt{compute\_objective(s, shocks, params, mu=100)}

Fungsi utama yang menghitung nilai $J_{penalized}(s)$ untuk sebuah state. Fungsi ini memanggil \texttt{simulate\_economy}, seluruh fungsi proksi, dan \texttt{check\_constraints} secara berurutan.

\begin{itemize}
    \item \texttt{s} (list[float]) state suku bunga.
    \item \texttt{shocks} (dict) guncangan stokastik.
    \item \texttt{params} (dict) seluruh parameter model, \textit{constraint}, dan bobot ($w_\pi, w_y, w_{pp}, w_r, \delta$, serta parameter yang dibutuhkan \texttt{simulate\_economy} dan \texttt{check\_constraints}).
    \item \texttt{mu} (float) bobot penalti \textit{constraint}, default 100.
    \item \textit{Return} \texttt{float} nilai $J_{penalized}(s) = -L(s) - \mu \sum_{i=2}^{5} \text{violation}_i(s)^2$.
\end{itemize}

\texttt{compute\_rmsd(trajectory, pi\_star=2.5)}

Menghitung RMSD$_{\pi}$ dari \textit{trajectory} hasil simulasi. Fungsi ini murni menghitung metrik evaluasi tanpa melibatkan bobot subjektif atau penalti \textit{constraint}.

\begin{itemize}
    \item \texttt{trajectory} (dict) hasil simulasi dari \texttt{simulate\_economy}, minimal mengandung kunci \texttt{"pi"}.
    \item \texttt{pi\_star} (float) target inflasi, default $2{,}5$.
    \item \textit{Return} \texttt{float} nilai RMSD$_{\pi} = \sqrt{\frac{1}{T} \sum (\pi_t - \pi^{*})^{2}}$.
\end{itemize}

\subsection*{Fase 3. Pencarian \textit{Local Search}}

\texttt{generate\_initial\_state(T=8, r\_min=3.0, r\_max=8.0, rng=None)}

Membentuk initial state $s_0$ secara acak dengan memilih setiap $r_t$ independen dari distribusi uniform pada $[r_{min}, r_{max}]$ lalu dibulatkan ke kelipatan 0,25\%, sehingga C1 terpenuhi secara implisit.

\begin{itemize}
    \item \texttt{T} (int) horizon kebijakan.
    \item \texttt{r\_min} (float) batas bawah suku bunga.
    \item \texttt{r\_max} (float) batas atas suku bunga.
    \item \texttt{rng} (numpy.random.Generator) generator bilangan acak.
    \item \textit{Return} \texttt{list[float]} state $s_0$ dengan panjang $T$.
\end{itemize}

\texttt{generate\_neighbor(s, move\_type=None, rng=None)}

Menghasilkan satu state tetangga dari $s$ menggunakan \textit{Move 1} atau \textit{Move 2} yang dipilih secara acak. Hasil modifikasi diperbaiki terhadap C1 sebelum dikembalikan.

\begin{itemize}
    \item \texttt{s} (list[float]) state asal.
    \item \texttt{move\_type} (int or None) jenis \textit{move}. Jika \texttt{None}, dipilih acak antara 1 dan 2.
    \item \texttt{rng} (numpy.random.Generator) generator bilangan acak.
    \item \textit{Return} \texttt{list[float]} state tetangga $s'$.
\end{itemize}

\texttt{generate\_neighbor\_crossover(s1, s2, rng=None)}

Menghasilkan dua state anak dari dua state induk melalui \textit{crossover} satu titik potong, khusus untuk \textit{Genetic Algorithm}.

\begin{itemize}
    \item \texttt{s1} (list[float]) state induk pertama.
    \item \texttt{s2} (list[float]) state induk kedua.
    \item \texttt{rng} (numpy.random.Generator) generator bilangan acak.
    \item \textit{Return} \texttt{tuple} dua \texttt{list[float]} sebagai state anak.
\end{itemize}

\texttt{hill\_climbing(s0, max\_iter, shocks, params)}

Algoritma \textit{Hill-Climbing} dasar. Mulai dari $s_0$, pada tiap iterasi membangkitkan satu tetangga. Jika $J(s') > J(s_{current})$, pindah ke $s'$. Jika tidak, tetap di $s_{current}$. Berhenti setelah \texttt{max\_iter} atau tidak ada perbaikan dalam $k$ iterasi berturut-turut.

\begin{itemize}
    \item \texttt{s0} (list[float]) initial state.
    \item \texttt{max\_iter} (int) jumlah maksimum iterasi.
    \item \texttt{shocks} (dict) guncangan stokastik.
    \item \texttt{params} (dict) parameter model dan bobot.
    \item \textit{Return} \texttt{dict} berisi \texttt{best\_state} (list[float]), \texttt{best\_score} (float), dan \texttt{history} (list[float]) nilai $J(s)$ per iterasi.
\end{itemize}

\texttt{simulated\_annealing(s0, max\_iter, T0, cooling\_rate, shocks, params)}

Algoritma \textit{Simulated Annealing}. Mirip \textit{Hill-Climbing} tetapi menerima tetangga lebih buruk dengan probabilitas $\exp((J(s') - J(s)) / T_{temp})$. Suhu $T_{temp}$ menurun secara geometrik tiap iterasi dengan faktor \texttt{cooling\_rate}.

\begin{itemize}
    \item \texttt{s0} (list[float]) initial state.
    \item \texttt{max\_iter} (int) jumlah maksimum iterasi.
    \item \texttt{T0} (float) suhu awal.
    \item \texttt{cooling\_rate} (float) faktor penurunan suhu per iterasi, $\in (0, 1)$.
    \item \texttt{shocks} (dict) guncangan stokastik.
    \item \texttt{params} (dict) parameter model dan bobot.
    \item \textit{Return} sama dengan \texttt{hill\_climbing}.
\end{itemize}

\texttt{genetic\_algorithm(pop\_size, generations, mutation\_rate, shocks, params)}

Algoritma \textit{Genetic Algorithm}. Diawali dengan \texttt{pop\_size} individu acak. Tiap generasi melakukan seleksi turnamen, \textit{crossover} satu titik, dan mutasi dengan probabilitas \texttt{mutation\_rate}. Elitisme menjaga individu terbaik tetap hidup ke generasi berikutnya.

\begin{itemize}
    \item \texttt{pop\_size} (int) jumlah individu dalam populasi.
    \item \texttt{generations} (int) jumlah generasi.
    \item \texttt{mutation\_rate} (float) probabilitas mutasi per individu, $\in (0, 1)$.
    \item \texttt{shocks} (dict) guncangan stokastik.
    \item \texttt{params} (dict) parameter model dan bobot.
    \item \textit{Return} \texttt{dict} berisi \texttt{best\_state} (list[float]), \texttt{best\_score} (float), dan \texttt{history} (list[float]) nilai $J(s)$ terbaik per generasi.
\end{itemize}

\subsection*{Fungsi Utilitas Tambahan}

\texttt{round\_to\_bps(value)}

Membulatkan nilai suku bunga ke kelipatan 25 basis poin terdekat.

\begin{itemize}
    \item \texttt{value} (float) nilai suku bunga dalam persen.
    \item \textit{Return} \texttt{float} nilai yang telah dibulatkan.
\end{itemize}

\texttt{clip\_state(s, r\_min=3.0, r\_max=8.0)}

Menjepit seluruh elemen state ke rentang $[r_{min}, r_{max}]$ dan membulatkannya. Berguna untuk memperbaiki state setelah modifikasi oleh \textit{move}.

\begin{itemize}
    \item \texttt{s} (list[float]) state yang akan diperbaiki.
    \item \texttt{r\_min} (float) batas bawah.
    \item \texttt{r\_max} (float) batas atas.
    \item \textit{Return} \texttt{list[float]} state yang telah memenuhi C1.
\end{itemize}

\section{Contoh Input dan Output}

Contoh berikut menggunakan seluruh parameter standar dari Bagian 14 spesifikasi. Kode ditulis dalam Python dengan dependensi \texttt{numpy} saja.

\subsection*{Contoh 1. Inisialisasi Parameter dan \textit{Shock}}

\begin{tcolorbox}[title=Input]
    \begin{lstlisting}[language=Python]
import numpy as np

T = 8
params = {
    "r0": 5.75, "pi0": 3.34, "y0": 0.00, "e0": 17885,
    "r0_US": 3.625, "pi_star": 2.5, "pp_star": 0.0,
    "rho_y": 0.60, "beta": 0.30, "kappa": 0.11, "phi": 0.30,
    "rho_pi": 0.85, "rho_e": 0.00, "r_star": 2.5,
    "alpha1": 0.20, "alpha2": 0.15, "rho_fiskal": 1.50,
    "r_min": 3.00, "r_max": 8.00,
    "delta_min": 1.50, "theta": 3.00, "sigma_max": 2.50,
    "delta": 0.95, "mu": 100,
    "w_pi": 1.0, "w_y": 0.5, "w_pp": 0.3, "w_r": 0.2,
}

shocks = generate_shocks(T=T, seed=42)
\end{lstlisting}
\end{tcolorbox}

\subsection*{Contoh 2. Pembangkitan Initial State}

Pemanggilan \texttt{generate\_initial\_state} menghasilkan state acak yang memenuhi C1 dan C2.

\begin{tcolorbox}[title=Input]
    \begin{lstlisting}[language=Python]
rng = np.random.default_rng(42)
s0 = generate_initial_state(T=8, r_min=3.0, r_max=8.0, rng=rng)
\end{lstlisting}
\end{tcolorbox}

\begin{tcolorbox}[title=Output]
    \begin{lstlisting}[language=Python]
>>> s0
[6.00, 6.25, 6.25, 6.00, 5.75, 5.50, 5.25, 5.25]
\end{lstlisting}
\end{tcolorbox}

\subsection*{Contoh 3. Simulasi Ekonomi}

State $s_0$ disimulasikan pada model makroekonomi untuk memperoleh \textit{trajectory} variabel.

\begin{tcolorbox}[title=Input]
    \begin{lstlisting}[language=Python]
traj = simulate_economy(s0, shocks, params)
\end{lstlisting}
\end{tcolorbox}

\begin{tcolorbox}[title=Output]
    \begin{lstlisting}[language=Python, basicstyle=\ttfamily\tiny]
>>> for t in range(T):
...     print(f"t={t+1}: r={s0[t]:.2f}%  y={traj['y'][t]:+.2f}  "
...           f"pi={traj['pi'][t]:.2f}%  e={traj['e'][t]:.0f}  "
...           f"PP={traj['PP'][t]:+.2f}  CA={traj['CA'][t]:+.2f}%")

t=1: r=6.00%  y=-0.12  pi=3.58%  e=17742  PP=-3.70  CA=-1.18%
t=2: r=6.25%  y=-0.35  pi=3.13%  e=17605  PP=-3.48  CA=-0.87%
t=3: r=6.25%  y=-0.55  pi=2.97%  e=17480  PP=-3.52  CA=-0.61%
t=4: r=6.00%  y=-0.48  pi=2.81%  e=17368  PP=-3.28  CA=-0.43%
t=5: r=5.75%  y=-0.30  pi=2.72%  e=17287  PP=-3.02  CA=-0.27%
t=6: r=5.50%  y=-0.12  pi=2.65%  e=17231  PP=-2.77  CA=-0.13%
t=7: r=5.25%  y=+0.04  pi=2.59%  e=17200  PP=-2.55  CA=-0.02%
t=8: r=5.25%  y=+0.13  pi=2.55%  e=17185  PP=-2.43  CA=+0.07%
\end{lstlisting}
\end{tcolorbox}

\subsection*{Contoh 4. Evaluasi Objective dan \textit{Constraint}}

\begin{tcolorbox}[title=Input]
    \begin{lstlisting}[language=Python]
result = check_constraints(s0, traj, params)
score = compute_objective(s0, shocks, params)
rmsd = compute_rmsd(traj, pi_star=2.5)
\end{lstlisting}
\end{tcolorbox}

\begin{tcolorbox}[title=Output]
    \begin{lstlisting}[language=Python]
>>> result["feasible"]
True

>>> result["violations"]
{'C2': 0.0, 'C3': 0.0, 'C4': 0.0}

>>> score
-8.935

>>> rmsd
0.476
\end{lstlisting}
\end{tcolorbox}

\subsection*{Contoh 5. Menjalankan Ketiga Algoritma}

\begin{tcolorbox}[title=Input]
    \begin{lstlisting}[language=Python]
rng = np.random.default_rng(42)
s0 = generate_initial_state(T=8, r_min=3.0, r_max=8.0, rng=rng)

result_hc = hill_climbing(s0, max_iter=1000, shocks=shocks,
                           params=params)

rng = np.random.default_rng(42)
sa_s0 = generate_initial_state(T=8, r_min=3.0, r_max=8.0, rng=rng)
result_sa = simulated_annealing(sa_s0, max_iter=1000, T0=10.0,
                                 cooling_rate=0.995, shocks=shocks,
                                 params=params)

result_ga = genetic_algorithm(pop_size=50, generations=200,
                               mutation_rate=0.15, shocks=shocks,
                               params=params)

rmsd_hc = compute_rmsd(simulate_economy(result_hc["best_state"],
                       shocks, params))
rmsd_sa = compute_rmsd(simulate_economy(result_sa["best_state"],
                       shocks, params))
rmsd_ga = compute_rmsd(simulate_economy(result_ga["best_state"],
                       shocks, params))
\end{lstlisting}
\end{tcolorbox}

\begin{tcolorbox}[title=Output]
    \begin{lstlisting}[language=Python]
>>> print(f"HC  : J={result_hc['best_score']:.4f}  "
...       f"state={result_hc['best_state']}")
HC  : J=-8.9352  state=[6.00, 6.25, 6.25, 6.00, 5.75, 5.50, 5.25, 5.25]

>>> print(f"SA  : J={result_sa['best_score']:.4f}  "
...       f"state={result_sa['best_state']}")
SA  : J=-5.1823  state=[5.75, 5.50, 5.25, 5.00, 5.00, 5.00, 4.75, 4.75]

>>> print(f"GA  : J={result_ga['best_score']:.4f}  "
...       f"state={result_ga['best_state']}")
GA  : J=-4.9571  state=[5.75, 5.50, 5.25, 5.00, 5.00, 4.75, 4.75, 4.75]
\end{lstlisting}
\end{tcolorbox}

\subsection*{Contoh 6. Konvergensi \textit{Loss} per Iterasi}

\begin{tcolorbox}[title=Output (nilai $J(s)$ pada beberapa iterasi terpilih)]
    \begin{lstlisting}[language=Python]
=== Hill-Climbing Convergence (max_iter=1000) ===
Iter    10: J(s) = -9.2431
Iter    50: J(s) = -9.0124
Iter   100: J(s) = -8.9978
Iter   200: J(s) = -8.9352
Iter   500: J(s) = -8.9352   (stagnant)
Iter  1000: J(s) = -8.9352   (best found)

=== Simulated Annealing Convergence (max_iter=1000) ===
T=10.000  Iter    10: J(s) = -9.0124  best=-9.0124
T= 6.083  Iter    50: J(s) = -7.8412  best=-7.8412
T= 2.509  Iter   200: J(s) = -5.9821  best=-5.9821
T= 0.897  Iter   500: J(s) = -5.1835  best=-5.1823
T= 0.033  Iter  1000: J(s) = -5.1826  best=-5.1823  (SA TERMINATED)
\end{lstlisting}
\end{tcolorbox}

\subsection*{Contoh 7. Perbandingan Akhir Tiga Algoritma}

\begin{tcolorbox}[title=Output (Ringkasan Perbandingan)]
    \begin{lstlisting}[language=Python]
===================================================================
           J(s)     RMSD_pi   #Eval   Waktu(s)  Feasible  pi_T
-------------------------------------------------------------------
HC        -8.9352    0.476      523      2.14      YES     2.55%
SA        -5.1823    0.241     1000      4.08      YES     2.48%
GA        -4.9571    0.218    10000*     5.73      YES     2.47%
-------------------------------------------------------------------
*GA: 50 individu x 200 generasi

Best J(s):        Genetic Algorithm (-4.9571)
Best RMSD_pi:     Genetic Algorithm (0.218)
Interpretasi:     0.218 < 0.50 -> SANGAT BAIK
\end{lstlisting}
\end{tcolorbox}

